\documentclass{article}
\usepackage{graphicx}
\usepackage{amsmath}
\usepackage[T2A]{fontenc}  % For Cyrillic fonts
\usepackage[utf8]{inputenc}  % For UTF-8 encoding
\usepackage[russian,english]{babel}  % For Russian language support
\usepackage{amssymb}
\usepackage[margin=2cm]{geometry}
\usepackage{listings}
\usepackage{xcolor}
\usepackage{float}
\usepackage{hyperref}

\lstset{
    language=Python,
    basicstyle=\ttfamily\small,
    keywordstyle=\color{blue},
    stringstyle=\color{red},
    commentstyle=\color{gray},
    backgroundcolor=\color{lightgray!20},
    frame=single,
    rulecolor=\color{black},
    numbers=left,
    numberstyle=\tiny\color{gray},
    stepnumber=1,
    numbersep=10pt,
    showspaces=false,
    showstringspaces=false,
    breaklines=true,
    breakatwhitespace=true,
    tabsize=4,
    captionpos=b
}

\title{Математическая статистика \\ Домашнее задание 2}
\author{Никита Загайнов}

\date{April 2025}

\begin{document}

\maketitle

\begin{enumerate}
	\item
	      \begin{enumerate}
		      \item Используя факт, что $\sum_{i=1}^{n} X_i - n\theta \sim \Gamma(1, n)$, запишем доверительный интервал на эту величину:
		            \begin{equation*}
			            \begin{aligned}
				             & P \left( \Gamma_{\frac{\alpha}{2}} (1, n) < \sum_{i=1}^{n} X_i - n\theta < \Gamma_{1 - \frac{\alpha}{2}} (1, n) \right) = 1 - \alpha                                         \\
				             & P \left( \Gamma_{\frac{\alpha}{2}} (1, n) - \sum_{i=1}^{n} X_i <  - n\theta < \Gamma_{1 - \frac{\alpha}{2}} (1, n) - \sum_{i=1}^{n} X_i \right) = 1 - \alpha                 \\
				             & P \left( \frac{\sum_{i=1}^{n} X_i - \Gamma_{1 - \frac{\alpha}{2}} (1, n)}{n} < \theta < \frac{\sum_{i=1}^{n} X_i - \Gamma_{\frac{\alpha}{2}} (1, n)}{n} \right) = 1 - \alpha \\
			            \end{aligned}
		            \end{equation*}

		      \item Используем ЦПТ:
		            \[
			            \sqrt{n} \frac{\sum_{i=1}^n X_i - n\mathbb{E}X}{\sigma(X)} = \sqrt{n} \left( \sum_{i=1}^n X_i - n(1 + \theta) \right) \sim \mathcal{N}(0, 1)
		            \]
		            для больших $n$. Для
		            построения доверительного интервала используем нормальное распределение:
		            \begin{equation*}
			            \begin{aligned}
				             & P \left( \Phi_{\frac{\alpha}{2}}^{-1} < \sqrt{n} \left( \sum_{i=1}^n X_i - n(1 + \theta) \right) < \Phi_{1 - \frac{\alpha}{2}}^{-1} \right) = 1 - \alpha                                                         \\
				             & P \left( \frac{\Phi_{\frac{\alpha}{2}}^{-1}}{\sqrt{n}} - \sum_{i=1}^n X_i < - n(1 + \theta) < \frac{\Phi_{1 - \frac{\alpha}{2}}^{-1}}{\sqrt{n}} - \sum_{i=1}^n X_i \right) = 1 - \alpha                          \\
				             & P \left( \frac{1}{n} \sum_{i=1}^n X_i - \frac{\Phi_{1 - \frac{\alpha}{2}}^{-1}}{n\sqrt{n}} - 1 < \theta < \frac{1}{n} \sum_{i=1}^n X_i - \frac{\Phi_{\frac{\alpha}{2}}^{-1}}{n\sqrt{n}} - 1 \right) = 1 - \alpha \\
			            \end{aligned}
		            \end{equation*}
	      \end{enumerate}

	\item
	      \begin{enumerate}
		      \item Используем величину $X_{(n)}$, распределённую следующим образом:
		            \[
			            F(x) = P \left( X_{(n)} < x \right) = \prod_{i=1}^{n} P \left( X_i < x \right) = \left( \frac{x}{\theta} \right)^n
		            \]

		            Запишем доверительный интервал для этой случайной величины:
		            \begin{equation*}
			            \begin{aligned}
				             & P \left( F^{-1}_{\frac{\alpha}{2}} < X_{(n)} <  F^{-1}_{1 - \frac{\alpha}{2}} \right) = 1 - \alpha                                                                                     \\
				             & P \left( \theta \left( \frac{\alpha}{2} \right)^{\frac{1}{n}} < X_{(n)} < \theta \left( 1 - \frac{\alpha}{2} \right)^{\frac{1}{n}} \right) = 1 - \alpha                                \\
				             & P \left( \frac{1}{X_{(n)}} \left( \frac{\alpha}{2} \right)^{\frac{1}{n}} < \frac{1}{\theta} < \frac{1}{X_{(n)}} \left( 1 - \frac{\alpha}{2} \right)^{\frac{1}{n}} \right) = 1 - \alpha \\
				             & P \left( X_{(n)} \frac{1}{\left( 1 - \frac{\alpha}{2} \right)^{\frac{1}{n}}} < \theta < X_{(n)} \frac{1}{\left( \frac{\alpha}{2} \right)^{\frac{1}{n}}} \right) = 1 - \alpha           \\
			            \end{aligned}
		            \end{equation*}

		      \item По ЦПТ:
		            \[
			            \sqrt{n} \frac{\sum_{i=1}^n X_i - n \frac{\theta}{2}}{\frac{1}{12} \theta^2} = 6 \sqrt{n} \frac{\sum_{i=1}^n X_i - n}{\theta} \sim \mathcal{N}(0, 1)
		            \]
		            Доверительный интервал:
		            \begin{equation*}
			            \begin{aligned}
				             & P \left( \Phi_{\frac{\alpha}{2}}^{-1} < 6 \sqrt{n} \frac{\sum_{i=1}^n X_i - n}{\theta} < \Phi_{1 - \frac{\alpha}{2}}^{-1} \right) = 1 - \alpha                                                                       \\
				             & P \left( \frac{6 \sqrt{n} \left( \sum_{i=1}^n X_i - n \right)}{\Phi_{1 - \frac{\alpha}{2}}^{-1}} < \theta < \frac{6 \sqrt{n} \left( \sum_{i=1}^n X_i - n \right)}{\Phi_{\frac{\alpha}{2}}^{-1}} \right) = 1 - \alpha \\
			            \end{aligned}
		            \end{equation*}
	      \end{enumerate}

	\item
	      Найдём критерий уровня значимости $\alpha$ вида $X_{(n)} \leq c \theta_0$:
	      \begin{equation*}
		      \begin{aligned}
			       & P_1 \left( X_{(n)} \leq c \theta_0 \right) = \alpha   \\
			       & \left( \frac{c \theta_0}{\theta_0} \right)^n = \alpha \\
			       & c = \alpha^{\frac{1}{n}}                              \\
		      \end{aligned}
	      \end{equation*}
	      Докажем, что этот критерий равномерно наиболее мощный. Отношение правдоподобий:
	      \[
		      \Lambda(X) = \frac{L(\theta)}{L(\theta_0)} = \frac{\theta_0^n I(X_{(n)} \leq \theta)}{\theta^n I(X_{(n)} \leq \theta_0)}
	      \]
	      Значение отношения монотонно возрастает при уменьшении $\theta$ (при $\theta > X_{(n)}$), что значит, что этот критерий наиболее мощный.

	\item
	      Запишем отношение правдоподобий:
	      \[
		      \Lambda(X) = \frac{L(\theta_1)}{L(\theta_0)} = \frac{\theta_1^n \exp(-\theta_1 (\sum_{i=1}^n X_i))}{\theta_0^n \exp(-\theta_0 (\sum_{i=1}^n X_i))} =\frac{L(\theta_1)}{L(\theta_0)} = \frac{\theta_1^n}{\theta_0^n} \exp \left( (\theta_0 - \theta_1) \sum_{i=1}^n X_i \right)
	      \]
	      По лемме Неймана-Пирсона, оптимальный критерий при заданном уровне значимости $\alpha$ имеет вид $P_0(\Lambda(X) > \lambda) = \alpha$, для некоторого $\lambda$. Подставив значение $\Lambda(X)$ в это неравенство и проведя некоторые преобразования, получим:
	      \[
		      P_0 \left( \sum_{i=1}^n X_i > \lambda \right) = \alpha
	      \]
	      где $\lambda = \Gamma^{-1}_{1 - \alpha}(\theta_0, n)$ - квантиль уровня $1 - \alpha$ распределения гамма-распределения (распределения суммы $n$ независимых экспоненциальных случайных величин с параметром $\theta_0$). Тогда РНМК будет иметь вид:
	      \[
		      \sum_{i=1}^n X_i > \Gamma^{-1}_{1 - \alpha}(\theta_0, n)
	      \]

	\item
	      Отношение правдоподобий:
	      \[
		      \Lambda(X) = \frac{e^{-n\theta_1} \theta_1^{\sum_{i=1}^n X_i}\prod_{i=1}^n \frac{1}{X_i!}}{e^{-n\theta_0} \theta_0^{\sum_{i=1}^n X_i}\prod_{i=1}^n \frac{1}{X_i!}} = \exp (n(\theta_0 - \theta_1)) \frac{\theta_1^{\sum_{i=1}^n X_i}}{\theta_0^{\sum_{i=1}^n X_i}}
	      \]
	      Применим лемму Неймана-Пирсона, применив несколько преобразований к полученному выражению:
	      \[
		      P_0 \left( \sum_{i=1}^n X_i > \lambda \right) = \alpha
	      \]
	      Сумма пуассоновских величин следует распределению $Poi(n \theta)$, тогда при заданном уровне значимости $\alpha$ уравнение сводится к
	      \[
		      P_0 \left( \sum_{i=1}^n X_i > Poi^{-1}_{1 - \alpha} (n \theta) \right) = \alpha
	      \]
	      РМНК имеет вид
	      \[
		      \sum_{i=1}^n X_i > Poi^{-1}_{1 - \alpha} (n \theta)
	      \]

\end{enumerate}


\end{document}
